\chapter{Pendahuluan \& Visi Luna AI}

\section{Latar Belakang}
Luna AI adalah platform konseling kesehatan mental berbasis Kecerdasan Buatan (\textit{Artificial Intelligence}) yang dirancang untuk menyediakan ruang aman, inklusif, dan responsif 24/7 bagi pengguna. Luna AI memadukan modalitas \textit{Natural Language Processing} (NLP) canggih, evaluasi risiko real-time melalui \textit{Safety Gates}, serta sintesis harian untuk mendukung kesejahteraan emosional pengguna.

\begin{lunacallout}{Visi Utama Luna AI}
Menjadi pendamping emosional terpercaya berbasis AI yang mampu memberikan pertolongan pertama kesehatan mental, mendukung refleksivitas emosional, dan mendeteksi potensi krisis emosional secara proaktif.
\end{lunacallout}

\section{Tujuan Dokumentasi Rule Book}
Dokumen ini disusun sebagai panduan resmi (\textit{Official Rule Book \& Operations Manual}) yang bertujuan untuk:
\begin{enumerate}[label=\alph*.]
    \item Memberikan panduan penggunaan antarmuka aplikasi Luna AI secara komprehensif.
    \item Menjelaskan aturan etika, batasan interaksi, dan norma konseling AI.
    \item Menguraikan protokol keselamatan (\textit{Safety Gates}) dan mekanisme penanganan situasi krisis emosional.
    \item Menyediakan rujukan visual melalui area tangkapan layar (\textit{screenshot placeholders}) untuk alur operasional utama.
\end{enumerate}

\section{Struktur Dokumen}
Rule Book ini terbagi menjadi 5 bab utama:
\begin{itemize}
    \item \textbf{Bab 1: Pendahuluan \& Visi Luna AI} — Pengenalan umum dan tujuan operasional.
    \item \textbf{Bab 2: Panduan Antarmuka \& Flow Aplikasi} — Alur penggunaan fitur beserta tangkapan layar.
    \item \textbf{Bab 3: Etika \& Aturan Konseling AI} — Aturan tata kelola konseling dan batasan AI.
    \item \textbf{Bab 4: Protokol Keselamatan \& Safety Gates} — Manajemen risiko emosional dan rujukan darurat.
    \item \textbf{Bab 5: Troubleshooting \& FAQ} — Solusi masalah teknis dan pertanyaan umum.
\end{itemize}
